% !PRINTED 648
\weberpage{648}
%<<CONT
conjugirt imaginär sind; denn erstens sind ihre $n^{\mathrm{ten}}$ Potenzen conjugirt
imaginär, sie selbst also zunächst bis auf eine $n^{\mathrm{te}}$ Einheitswurzel als
Factor. Da aber ihr Product reell ist, so muss diese $n^{\mathrm{te}}$ Einheitswurzel
$= 1$ sein. Da nun ausserdem
\[
r_\nu \equiv - r_{\nu + \frac{n-1}{2}} \pmod{n}
\]
ist, so zeigt die Formel (11), dass in diesem Falle die Wurzeln $\xi_z$ alle reell sind.

Damit ist also der Satz für einen reellen Körper $\mathfrak{K}$ bewiesen:

\gesperrt{Hat die cyklische Gleichung, deren Wurzeln die $k_0$, $k_1 \ldots k_{n-2}$ sind, reelle
Wurzeln, so ist von den Wurzeln $\xi$ eine reell, die übrigen imaginär; hat diese
cyklische Gleichung imaginäre Wurzeln, so sind alle $\xi$ reell.}

\begin{center}
§. 188.

\gesperrt{Metacyklische Gleichungen fünften Grades.}
\end{center}

Durch die Sätze des vorangehenden Paragraphen sind die Wurzeln einer metacyklischen
Gleichung von Primzahlgrad in eine allgemein gültige Form gebracht, die es gestattet,
alle diese Grössen wirklich zu bilden, wenn man die Kenntniss der Wurzeln cyklischer
Gleichungen voraussetzt. Ganze Systeme cyklischer Gleichungen jeden Grades liefert uns
z. B. die Kreistheilungstheorie, deren Wurzeln die Kreistheilungsperioden sind. So
können wir also beispielsweise für den Körper der rationalen Zahlen beliebig viele
metacyklische Gleichungen bilden. Dass darin alle diese Gleichungen im Körper der
rationalen Zahlen enthalten sind, ist ein sehr merkwürdiger von \gesperrt{Kronecker}
herrührender Satz, den wir im zweiten Bande kennen lernen werden. Hier wollen wir als
Anwendung und Veranschaulichung des Vorhergehenden noch die specielle Aufgabe
behandeln, in einem beliebigen Körper $\mathfrak{K}$ die Wurzeln aller metacyklischen
Gleichungen fünften Grades zu finden.

Dazu ist erforderlich und hinreichend, dass wir die Wurzeln $k_0, k_1, k_2, k_3$ einer
cyklischen Gleichung vierten Grades allgemein bestimmen, und zwar unter der
Voraussetzung, dass \gesperrt{$k_0, k_1, k_2, k_3$ unter einander verschieden und keine
von ihnen Null sei}.
